% !TEX root = ../main.tex
% =====================================================================
% 6. MY PROPOSED MIL POOLING HEADS          target: 4 - 4.5 pages
%
% The mathematical core of the report. The notation deliberately follows
% the MILLET paper (bag / instance / attention / bag logits), so a reader
% who knows that paper can compare my heads with theirs line by line.
%
% Rules I follow in this section:
%   - one head = one short motivation + one equation block + one property
%   - every head must reduce to the Conjunctive baseline for some value
%     of its parameters, and I state that as a Property
%   - no results here; results are in section 8
% =====================================================================

\section{Proposed MIL Pooling Heads}
\label{sec:pooling}

\begin{contribution}
The seven pooling heads in this section were designed and implemented by me. Five of
them (\Cref{sec:pool:classwise}--\Cref{sec:pool:attnmax}) come from my own analysis of
the three weaknesses of Conjunctive pooling listed in \Cref{sec:bg:millet}. The last
two (\Cref{sec:pool:gated}--\Cref{sec:pool:dsmil}) adapt ideas published for
whole-slide image classification \citep{ilse2018attention, li2021dsmil} to the time
series setting; the adaptation, the per-class formulation and the implementation are
mine.
\end{contribution}

\begin{guide}
Order the subsections so that they mirror the weakness list of \Cref{sec:bg:millet}
(class-agnostic gate $\rightarrow$ unnormalised weights $\rightarrow$ fixed mean).
Before the maths, always give one plain sentence: ``the problem is X, the fix is Y''.
\textbf{Questions to answer for every head:} What can it express that Conjunctive
cannot? What does it cost (extra parameters, extra compute)? What did I expect it to
improve -- accuracy, AOPCR, or both?
\end{guide}

% ---------------------------------------------------------------------
\subsection{Notation and the common template}
\label{sec:pool:notation}

\begin{table}[h]
\caption{Notation used throughout this section. It follows \citet{early2024millet} so
that my heads and theirs can be read side by side.}
\label{tab:notation}
\begin{center}
\small
\begin{tabular}{ll}
\toprule
\textbf{Symbol} & \textbf{Meaning} \\
\midrule
$\bag = (\inst_1,\dots,\inst_{\nsteps})$ & the bag: one time series of length $\nsteps$ \\
$j \in \{1,\dots,\nsteps\}$              & instance index (one time step) \\
$k \in \{1,\dots,\nclz\}$                & class index \\
$\emb_j \in \R^{\dmodel}$                & encoder embedding of time step $j$ \\
$\instpred_j \in \R^{\nclz}$             & per-time-step class logits \\
$a_j$ or $a_j^k$                         & attention gate (shared, or one per class) \\
$g_j^k$                                  & gated evidence $a_j^k \instpred_j^k$ \\
$\bagpred \in \R^{\nclz}$                & bag logits (the prediction for the series) \\
$\interp \in \R^{\nclz \times \nsteps}$  & interpretation: importance of step $j$ for class $k$ \\
$\dattn$                                 & width of the attention hidden layer \\
\bottomrule
\end{tabular}
\end{center}
\end{table}

Every head I propose keeps the same two-part template as Conjunctive pooling: an
attention branch that decides \emph{where} to look, and a per-time-step classifier
that decides \emph{what} each place argues for,
\begin{equation}
  \instpred_j = W_c \emb_j + \mathbf{b}_c \in \R^{\nclz},
  \qquad
  \mathbf{a}_j = \mathcal{A}(\emb_j) ,
  \label{eq:common_template}
\end{equation}
and they differ only in the shape of $\mathcal{A}$ and in how the $\nsteps$
per-time-step contributions are aggregated into $\bagpred$.

\begin{guide}
\todo{Optional but strong: add one short paragraph explaining why the interpretation
$\interp$ must be exactly the quantity the bag logit is built from. That is the
property that makes the explanation faithful, and it is the reason I never compute
$\interp$ with a separate branch.}
\end{guide}
{This is worth stating explicitly, because it is the reason the explanation can be
trusted: for every head in this section, $\interp$ is not computed by a separate branch
after the fact -- it is exactly the per-time-step, per-class quantity that $\bagpred$ is
built by summing or averaging (for example, \Cref{eq:cw_bag} sets $\interp_{k,j}=g_j^k$,
the very term summed to build $\bagpred^{\,k}$). This matches what
\texttt{seanet/pooling.py} actually returns: every \texttt{forward()} method builds
\texttt{bag\_logits} and \texttt{interpretation} from the same intermediate tensor
(\texttt{g} or \texttt{contrib}), never from two independent computations. If the two
were computed separately, nothing would stop them from disagreeing, and the explanation
would lose exactly the property \Cref{sec:bg:related} says post-hoc explainers lack.}

% ---------------------------------------------------------------------
\subsection{Class-wise conjunctive pooling}
\label{sec:pool:classwise}

\textbf{Problem.} \todo{One sentence: the gate in \Cref{eq:conj_gate} is a single
number per time step, shared by all classes.}
{Because $a_j$ in \Cref{eq:conj_gate} is one number shared by every class, a time step
that supports class A cannot, at the same time, argue against class B -- one gate value
has to serve all $\nclz$ classes at once.}

\textbf{Fix.} One gate per class, produced by the same small attention MLP with
$\nclz$ outputs instead of one:
\begin{align}
  \mathbf{a}_j &= \sig\!\left(W_2 \tanh(W_1 \emb_j)\right) \in [0,1]^{\nclz},
  \qquad W_1 \in \R^{\dattn \times \dmodel},\; W_2 \in \R^{\nclz \times \dattn},
  \label{eq:cw_gate}\\
  g_j^k &= a_j^k \, \instpred_j^{\,k},
  \label{eq:cw_evidence}\\
  \bagpred^{\,k} &= \frac{1}{\nsteps}\sum_{j=1}^{\nsteps} g_j^k,
  \qquad
  \interp_{k,j} = g_j^k .
  \label{eq:cw_bag}
\end{align}

\begin{property}[Strict generalisation]
\label{prop:cw}
If $a_j^k = a_j$ for every class $k$, \Cref{eq:cw_bag} reduces exactly to Conjunctive
pooling \Cref{eq:conj_bag}; if $a_j^k \to 1$ it reduces to instance pooling.
\end{property}
\todo{One sentence on why this matters: the head can only \emph{add} capacity on top
of the baseline, so a well-trained model should never do worse.}
{This matters because \Cref{prop:cw} means the head only ever \emph{adds} capacity: at
worst it can copy the Conjunctive baseline exactly (by collapsing every $a_j^k$ to the
same value), so a well-trained model has no reason to score below it.}

% ---------------------------------------------------------------------
\subsection{Softmax conjunctive pooling with learnable temperature}
\label{sec:pool:softmax}

\textbf{Problem.} \todo{One sentence: sigmoid gates are unnormalised, and dividing by
$\nsteps$ makes the bag score depend on the length of the series.}
{Dividing the sigmoid gate's sum by $\nsteps$ in \Cref{eq:conj_bag} means the bag score's
scale shifts with the length of the series: two series with identical evidence but
different lengths $\nsteps$ do not get the same score, which is a problem across
datasets whose series length varies.}

\textbf{Fix.} Turn the attention into a real probability distribution over time, with
a learnable temperature $\tau$:
\begin{align}
  s_j^k     &= \left(W_2 \tanh(W_1 \emb_j)\right)_k \in \R,
  \label{eq:sm_score}\\
  \alpha_j^k &= \softmaxt\!\left(\frac{s_j^k}{\tau}\right)
              = \frac{\exp(s_j^k/\tau)}{\sum_{i=1}^{\nsteps}\exp(s_i^k/\tau)},
  \qquad \sum_{j=1}^{\nsteps}\alpha_j^k = 1,
  \label{eq:sm_alpha}\\
  \bagpred^{\,k} &= \sum_{j=1}^{\nsteps} \alpha_j^k \, \instpred_j^{\,k},
  \qquad
  \interp_{k,j} = \alpha_j^k \, \instpred_j^{\,k}.
  \label{eq:sm_bag}
\end{align}
The temperature is stored as $\tau = \exp(\rho)$ with $\rho$ learnable, so it stays
strictly positive during training.

\begin{property}[Interpolation between mean and max]
\label{prop:sm}
As $\tau \to \infty$, $\alpha_j^k \to 1/\nsteps$ and \Cref{eq:sm_bag} becomes a plain
mean over time; as $\tau \to 0$ it selects the single highest-scoring time step.
Because the weights sum to one, the bag score does not depend on $\nsteps$.
\end{property}

% ---------------------------------------------------------------------
\subsection{Adaptive class-wise pooling}
\label{sec:pool:adaptive}

\textbf{Problem.} \todo{One sentence: some datasets hide the evidence in one short
event, others spread it over the whole series; a fixed mean cannot serve both.}
{On some datasets the evidence is a single short event -- a spike that a max-like
aggregator should highlight -- while on others it is spread across most of the series,
where an average is the safer choice; the fixed mean in \Cref{eq:conj_bag} can only ever
be right for one of these two cases.}

\textbf{Fix.} Keep the per-class gates of \Cref{eq:cw_gate}, then aggregate with a
learnable sharpness $\beta_k \geq 0$ per class (a numerically stable soft-argmax):
\begin{align}
  g_j^k &= a_j^k \, \instpred_j^{\,k},
  \qquad \beta_k = \softplus(\theta_k) \geq 0,
  \label{eq:ad_beta}\\
  w_j^k &= \softmaxt\!\left(\beta_k \, g_j^k\right),
  \label{eq:ad_weights}\\
  \bagpred^{\,k} &= \sum_{j=1}^{\nsteps} w_j^k \, g_j^k,
  \qquad
  \interp_{k,j} = g_j^k .
  \label{eq:ad_bag}
\end{align}

\begin{property}[Learned mean--max aggregator]
\label{prop:ad}
$\beta_k \to 0$ gives $w_j^k \to 1/\nsteps$, i.e.\ class-wise conjunctive pooling;
$\beta_k \to \infty$ gives $\bagpred^{\,k} \to \max_j g_j^k$. Since $\bagpred^{\,k}$
is always a convex combination of the values $g_j^k$, it stays inside their range,
which keeps training stable.
\end{property}
\todo{One sentence: $\beta_k$ starts small so training begins near the safe mean, and
the model sharpens it only where that helps.}
{$\beta_k$ starts small (\texttt{init\_beta} in \texttt{configs/models/sv2/seanet.yaml},
default $0.3$) so training begins close to the safe class-wise mean, and the optimiser
sharpens $\beta_k$ towards a max-like aggregator only for the classes where that
actually reduces the loss.}

% ---------------------------------------------------------------------
\subsection{Top-$k$ conjunctive pooling}
\label{sec:pool:topk}

\textbf{Problem.} \todo{One sentence: on a long series most time steps are background,
and averaging over all of them dilutes a handful of decisive points.}
{On a series with hundreds of time steps, most of them are just background: averaging
\Cref{eq:cw_bag} over all $\nsteps$ of them lets a handful of genuinely decisive points
get diluted by everything around them.}

\textbf{Fix.} Average only the $k$ strongest contributions of each class, with
$k$ a fixed fraction $\kappa$ of the length:
\begin{align}
  k &= \max\big(1, \lceil \kappa \nsteps \rceil\big), \qquad \kappa \in (0,1],
  \label{eq:topk_k}\\
  \mathcal{S}^k_{\text{top}} &= \big\{\, j : g_j^k \text{ is among the } k
    \text{ largest of } (g_1^k,\dots,g_{\nsteps}^k) \,\big\},
  \label{eq:topk_set}\\
  \bagpred^{\,k} &= \frac{1}{k}\sum_{j \in \mathcal{S}^k_{\text{top}}} g_j^k,
  \qquad
  \interp_{k,j} = g_j^k .
  \label{eq:topk_bag}
\end{align}

\begin{property}[Strict generalisation]
\label{prop:topk}
$\kappa = 1$ recovers class-wise conjunctive pooling exactly. For $\kappa < 1$ the
gradient reaches only the selected time steps, which concentrates the evidence.
\end{property}
\todo{One sentence on the expected effect on AOPCR: if the decision really rests on
few time steps, masking them should hurt the model more, so AOPCR should rise.}
{If the decision genuinely rests on the $k$ selected time steps, masking exactly those
points should hurt the model's confidence far more than masking $k$ random ones, so this
head is expected to raise AOPCR specifically, more than it is expected to raise raw
accuracy.}

% ---------------------------------------------------------------------
\subsection{Attention-max pooling}
\label{sec:pool:attnmax}

\textbf{Problem.} \todo{One sentence: same problem as above, but I wanted a simpler,
``hard'' version to check whether the smooth soft-argmax is really needed.}
{\Cref{sec:pool:adaptive} already fixes the fixed-mean problem, but its soft-argmax
mixes every time step a little; here I wanted the simplest possible ``hard''
alternative -- an explicit mean/max blend -- to check whether the smoother, learned
aggregator of \Cref{sec:pool:adaptive} is actually earning its extra complexity.}

\textbf{Fix.} A learnable per-class blend of the mean and the max:
\begin{equation}
  \bagpred^{\,k} = (1-\lambda_k)\,\frac{1}{\nsteps}\sum_{j=1}^{\nsteps} g_j^k
                   \; + \; \lambda_k \max_{j} g_j^k ,
  \qquad \lambda_k = \sig(\theta_k) \in [0,1].
  \label{eq:attnmax}
\end{equation}

\begin{property}
\label{prop:attnmax}
$\lambda_k \to 0$ recovers class-wise conjunctive pooling; $\lambda_k \to 1$ gives max
pooling. Each class learns on its own whether its evidence is a trend or a spike.
\end{property}
\todo{One sentence: this head exists to isolate one factor -- soft versus hard
selection -- which is what makes the comparison with
\Cref{sec:pool:adaptive} informative.}
{This head exists to isolate one factor: soft (weighted-average) versus hard (arg-max)
selection of the evidence. Because \Cref{eq:attnmax} and \Cref{eq:ad_bag} share the same
per-class gate and differ only in how they turn it into one number, comparing the two in
\Cref{sec:results} says whether a smooth soft-argmax is worth its extra complexity over
a plain, learned mean/max blend.}

% ---------------------------------------------------------------------
\subsection{Per-class gated attention}
\label{sec:pool:gated}

\begin{priorwork}
The gated attention mechanism itself is due to \citet{ilse2018attention} and is used
in the multi-scale MIL pipeline of \citet{hashimoto2020msdamil} for cancer subtype
classification. In those papers a bag is a whole-slide image and an instance is an
image patch.
\end{priorwork}

\begin{contribution}
\todo{One sentence: I transferred it to time series (bag = series, instance = time
step) and made the attention \emph{per class}, which the original formulation is not,
so that it produces a class-specific interpretation like my other heads.}
\end{contribution}

\begin{align}
  \mathbf{A}_j &= \tanh(V \emb_j) \odot \sig(U \emb_j) \in \R^{\dattn},
  \qquad V, U \in \R^{\dattn \times \dmodel},
  \label{eq:ga_features}\\
  s_j^k &= (W \mathbf{A}_j)_k, \qquad W \in \R^{\nclz \times \dattn},
  \label{eq:ga_score}\\
  \alpha_j^k &= \softmaxt(s_j^k),
  \qquad
  \bagpred^{\,k} = \sum_{j=1}^{\nsteps} \alpha_j^k \, \instpred_j^{\,k},
  \qquad
  \interp_{k,j} = \alpha_j^k \, \instpred_j^{\,k}.
  \label{eq:ga_bag}
\end{align}
\todo{Two sentences in words: the $\tanh$ branch proposes attention features and the
$\sig$ branch decides which of them to let through, so the gate is more selective than
$\tanh$ alone.}
{In words, the $\tanh$ branch of \Cref{eq:ga_features} proposes a set of candidate
attention features, while the $\sig$ branch looks at the same embedding and decides
which of those features to let through -- multiplying the two means a feature is only
used where the sigmoid branch says it is trustworthy. This makes the gate more selective
than a plain $\tanh$-only branch (\Cref{eq:conj_gate}), at the cost of one extra
$\dmodel \to \dattn$ layer (\texttt{seanet/pooling.py}, \texttt{GatedAttentionPooling},
the \texttt{attn\_gate} branch $U$). Unlike every other head in this section, this one
has no blend parameter that can be driven to zero, so it does \emph{not} reduce exactly
to Conjunctive pooling for any setting of its weights -- see \Cref{tab:pooling_summary}.}

% ---------------------------------------------------------------------
\subsection{Dual-stream conjunctive pooling}
\label{sec:pool:dsmil}

\begin{priorwork}
The two-stream idea is due to \citet{li2021dsmil} (DSMIL), again for whole-slide
images: one stream aggregates everything, a second stream focuses on the single most
critical instance and on the instances that resemble it.
\end{priorwork}

\begin{contribution}
\todo{One or two sentences: I dropped their contrastive pre-training, made the
critical instance and the query per class, and wrapped the two streams in a learnable
blend that \emph{starts} at the Conjunctive baseline, so the head can only match or
improve on it. My goal with this head was specifically the AOPCR gap.}
\end{contribution}

\textbf{Stream A (safe baseline).} The class-wise conjunctive mean of
\Cref{eq:cw_bag}:
\begin{equation}
  \bagpred^{\,k}_{\text{mean}} = \frac{1}{\nsteps}\sum_{j=1}^{\nsteps} g_j^k .
  \label{eq:ds_streamA}
\end{equation}

\textbf{Stream B (critical time step).} For each class, find the most supporting time
step, embed a small query for every time step, and re-weight by similarity to that
critical query:
\begin{align}
  m_k &= \argmax_{j} \; g_j^k ,
  \qquad
  \mathbf{q}_j = W_q \emb_j \in \R^{\dattn},
  \label{eq:ds_query}\\
  U_j^k &= \softmaxt\!\left(\langle \mathbf{q}_j, \mathbf{q}_{m_k}\rangle\right),
  \qquad
  \bagpred^{\,k}_{\text{ds}} = \sum_{j=1}^{\nsteps} U_j^k \, \instpred_j^{\,k}.
  \label{eq:ds_streamB}
\end{align}

\textbf{Blend.} With one learnable scalar $\lambda = \sig(\theta)$,
\begin{equation}
  \bagpred^{\,k} = (1-\lambda)\,\bagpred^{\,k}_{\text{mean}}
                 + \lambda\,\bagpred^{\,k}_{\text{ds}},
  \qquad
  \interp_{k,j} = (1-\lambda)\, g_j^k + \lambda\, U_j^k \instpred_j^{\,k}.
  \label{eq:ds_blend}
\end{equation}

\begin{property}[Safe by construction]
\label{prop:ds}
$\lambda \to 0$ recovers class-wise conjunctive pooling. The interpretation in
\Cref{eq:ds_blend} uses the same blend as the logit, so the explanation always matches
the decision. The attention is measured only against the single critical time step, so
the cost is $\mathcal{O}(\nsteps \dattn \nclz)$ and not $\mathcal{O}(\nsteps^2)$.
\end{property}

% ----- THE PROPOSED POOLING DIAGRAM GOES HERE ------------------------
% ----- DIAGRAM BRIEF for fig:pooling_arch (draw this in your own tool) -------------------
% Source of truth: seanet/pooling.py (ClasswiseConjunctivePooling, SoftmaxConjunctivePooling,
% AdaptiveClasswisePooling, TopKConjunctivePooling, AttentionMaxPooling, GatedAttentionPooling,
% DualStreamConjunctivePooling) and Eq. (common template) in Section 6.1.
%
% LEFT side (shared by every head), top to bottom:
%   1. Per-time-step embeddings z_j          (B, T, d)   <- input to every head below
%   2. TWO branches read z_j in parallel:
%        (a) attention branch A(z_j)          -> one gate value per time step [shape varies by head]
%        (b) classifier  p_j = W_c z_j + b_c   -> (B, T, C)   [IDENTICAL in every head]
%   3. Aggregation box: combines (a) and (b) into bag_logits (B, C) AND interpretation (B, C, T)
%      -- this box is the ONLY thing that changes between heads.
%
% RIGHT side: seven small boxes, one per head, each showing ONLY its aggregation rule (the
% shared left-hand side above is implied). One-line label + equation to put under each box:
%   - Conjunctive (baseline; draw in a different colour/outline -- reused, not mine)
%       a_j scalar, shared by all classes -> mean over time            [Eq. (conj_bag), Sec. 4.3]
%   - Class-wise            a_j^k per class -> mean over time           [Eq. (cw_bag)]
%   - Softmax + temperature a_j^k -> softmax over time / tau            [Eq. (sm_bag)]
%   - Adaptive class-wise   a_j^k -> soft-argmax with learned beta_k    [Eq. (ad_bag)]
%   - Top-k                 a_j^k -> mean of the top k over time        [Eq. (topk_bag)]
%   - Attention-max         a_j^k -> learned blend of mean and max      [Eq. (attnmax)]
%   - Gated attention       tanh*sigmoid features -> softmax over time  [Eq. (ga_bag)]
%   - Dual-stream           two streams blended, lambda near 0 at init  [Eq. (ds_blend)]
%
% Draw a short arrow from the Conjunctive box to EVERY other box EXCEPT Gated attention --
% it is the one head that does NOT reduce back to Conjunctive for any parameter setting
% (see Section 6.7's closing paragraph and Table 4).
%
% Suggested colours: Conjunctive = grey (reused, not mine); Class-wise / Softmax / Adaptive /
% Top-k / Attention-max = your accent colour (my own analysis, Sec. 6.2-6.6); Gated attention
% and Dual-stream = a second accent (adapted from the WSI MIL papers, Sec. 6.7-6.8).
% -------------------------------------------------------------------------------------------
\begin{figure}[t]
  \centering
  \figph{\textbf{Figure X -- Proposed pooling architecture.} One picture with the
  shared template on the left (embeddings $\rightarrow$ attention branch +
  per-time-step classifier) and the seven aggregation variants on the right, each
  drawn as a small box: shared gate (baseline), per-class gate, softmax over time,
  soft-argmax with $\beta$, top-$k$ selection, mean/max blend, gated attention,
  dual-stream. Arrows show where each one departs from the baseline.}
  \caption{\todo{Caption: all heads share the same skeleton; only the aggregation
  changes. Say which box is the \millet{} baseline and which boxes are mine.} The
  Conjunctive box (grey) is the \millet{} baseline; every other box is mine. Five
  (\Cref{sec:pool:classwise}--\Cref{sec:pool:attnmax}) share the left-hand template
  exactly and only change how the per-time-step evidence $g_j^k$ is aggregated over
  time; the last two (\Cref{sec:pool:gated}--\Cref{sec:pool:dsmil}) additionally change
  how the attention itself is computed, adapting ideas from whole-slide-image MIL.}
  \label{fig:pooling_arch}
\end{figure}

% ---------------------------------------------------------------------
\subsection{Summary and cost}
\label{sec:pool:summary}

\begin{table}[t]
\caption{My pooling heads at a glance. Every head reduces to the Conjunctive baseline
for the parameter value in the last column, so none of them can lose capacity.
\todo{Check the extra-parameter column against the code.}}
\label{tab:pooling_summary}
\begin{center}
\small
\begin{tabular}{p{0.20\linewidth} p{0.26\linewidth} p{0.24\linewidth} p{0.18\linewidth}}
\toprule
\textbf{Head} & \textbf{Weakness it fixes} & \textbf{Extra parameters} &
\textbf{Reduces to baseline when} \\
\midrule
Conjunctive \reused    & --                       & --    & -- \\
Class-wise             & shared gate              & $(\nclz-1)(\dattn+1)$ & $a_j^k = a_j$ \\
Softmax + temperature  & unnormalised weights     & $(\nclz-1)(\dattn+1) + 1$ & $\tau \to \infty$ \\
Adaptive class-wise    & fixed mean               & $(\nclz-1)(\dattn+1) + \nclz$ & $\beta_k \to 0$ \\
Top-$k$                & evidence dilution        & $(\nclz-1)(\dattn+1)$ & $\kappa = 1$ \\
Attention-max          & evidence dilution (hard) & $(\nclz-1)(\dattn+1) + \nclz$ & $\lambda_k \to 0$ \\
Gated attention        & weak gate                & $\dattn\dmodel + \dattn\nclz + \nclz - 1$ & \todo{--} no exact reduction (see \Cref{sec:pool:gated}) \\
Dual-stream            & scattered evidence       & $(\nclz-1)(\dattn+1) + \dattn\dmodel + \dattn + 1$ & $\lambda \to 0$ \\
\bottomrule
\end{tabular}
\end{center}
% Extra-parameter formulas derived directly from seanet/pooling.py: every head except Gated
% attention and Dual-stream reuses the SAME small attention MLP shape as ClasswiseConjunctivePooling
% (Linear(d_in,d_attn) -> Tanh -> Linear(d_attn,n_clz)), so their extra cost over the Conjunctive
% baseline's single-output gate is just the wider last layer, (n_clz-1)(d_attn+1), plus whatever
% extra learnable scalar(s) the head adds (temperature: +1; per-class beta or lambda: +n_clz;
% top-k's top_frac is a fixed hyperparameter, not a learned weight, so it adds 0). Gated attention
% is the one head that is NOT just a wider output layer -- it adds a whole second d_in->d_attn
% branch (attn_gate, the sigmoid gate U), which is why its extra-parameter cost has a d_attn*d_in
% term the others do not. Dual-stream adds ClasswiseConjunctivePooling's usual extra PLUS one more
% d_in->d_attn query layer (W_q) and one scalar blend. Sanity-check these against
% `python main.py params` before printing exact numbers for a specific d_attn/n_clz.
\end{table}

\begin{guide}
\textbf{TODO -- extra mathematics that could be added later, if space allows:}
\begin{itemize}\itemsep2pt
  \item A short proof of each Property (they are two lines each).
  \item The gradient of \Cref{eq:ad_bag} with respect to $\beta_k$, to explain why the
        head starts near the mean and sharpens only when it pays off.
  \item A note on the complexity of each head, in terms of $\nsteps$, $\dmodel$,
        $\dattn$ and $\nclz$.
  \item The attention-entropy penalty added to the training loss, written as an
        equation, and its weight $\lambda_{\text{focus}}$.
\end{itemize}
\end{guide}

\todo{Optional: write the training objective here, cross-entropy plus the attention
focus penalty, in one equation.}
{Every model in this report is trained with the same loss: label-smoothed
cross-entropy on the class prediction, plus a small penalty that rewards a
\emph{focused} attention (one that puts its weight on a few time steps rather than
spreading it over the whole series). The penalty is the entropy of the attention
returned by the pooling head (every head in this section exposes one, averaged over
classes into the $\bar a_j$ below); a peaked distribution has low entropy and a flat
one has high entropy, so minimising the loss pushes the attention to be peaked
(\texttt{seanet/train.py}, \texttt{attention\_entropy}).}
\begin{equation}
  \mathcal{L} = \underbrace{\mathrm{CE}(\bagpred, y)}_{\text{classification}}
              + \lambda_{\text{focus}} \cdot
                \underbrace{\Big({-}\sum_{j=1}^{\nsteps} \bar a_j \log(\bar a_j + \epsilon)\Big)}_{\text{attention entropy}} .
  \label{eq:loss}
\end{equation}
{Here $\bar a_j$ is the ``attn'' tensor each pooling head returns (the attention
averaged over classes, e.g.\ $\bar a_j = \frac{1}{\nclz}\sum_k a_j^k$ for the class-wise
family), $\epsilon$ avoids $\log 0$, and $\lambda_{\text{focus}}$ is \texttt{lambda\_entropy}
in the config -- default $0.01$, set to $0$ for the \millet{} baseline since the focus
penalty is a \seanet{} idea, not part of \millet{} (\texttt{configs/models/sv2/seanet.yaml},
\texttt{configs/models/sv1/millet.yaml}).}